\section{怎样又快又准地查字典}

我们学习语文，必须从认识汉字做起。为了确切地了解字的读音、写法、含义和用法，就必须学会查字典。

字典是按一定方式把字编排起来，供人查阅的工具书。它以单字为收集的对象，主要讲明单个字的念法、写法、意义和用法。由于汉字数以万计，实在难以一一记清，而那些认不得的字或弄不清的字，则往往成了我们学习语文的拦路虎，这就需要字典这位不说话的老师来帮助我们扫除障碍了。

查字典既要快又要准。光快不准，结果徒劳无功，还得重新查起；光准不快，花去很长时间，使人心烦意躁。那么，怎样才能做到又快又准呢？这除了常查字典以求熟能生巧之外，主要依靠掌握各种字典的检字法，熟悉各种字典的检字凡例。由于一般词典首先解说“字头”，就是先讲单字，然后在单字下列举词语，分词条注释意义，因此，我们掌握了字典的查法，也就可以掌握了词典的查法。在这儿，我们着重介绍几种字典检字法，掌握了检字法，经过一番实践，
\newpage

\noindent
自能做到又快又准地查字典。

1. 音序检字法

目前出版的多数字典或词典，如《新华字典》、《汉语词典》、《现代汉语词典》、《农民词典》、《汉语成语小词典》等，都是用这种检字法。

这些字典或词典，都是按《汉语拼音方案》字母表顺序排列，即按照二十六个拼音字母的顺序排列。由于《汉语拼音方案》规定，i、u、ü开头的韵母，自成音节时与前面有声母时写法不一样，就是说，这时i读为y，u换为w，ü换为yu；另外，v这个字母，只用来拼写外来语、少数民族语言和方言，而在普通话语音里不用，因此，一般按拼音字母顺序来排列的字典或词典，只用了二十三个拼音字母作为音序的标志，而没有i、u、v。为了查字方便，我们可以将字典或词典中用这二十三个拼音字母开头的音节，分别用标签标明，而将它们隔开，这样查起来就快当多了。

查字的时候，先要看所查的那个汉字字音的第一个字母是什么。例如查“稼”，字音的第一个字母是j，就立即翻到字典中用标签标明的“j”字母起头的部分，然后再根据拼音字母顺序，查看字音的第二个字母“i”，在“i”之后很快就能找到字音的第三个字母“a”。字母完全相同的字，再按阴、阳、上、去四种声调顺序去查，那么“jià”这个字音下面就出现了要查的“稼”字。由此可见，用音序检字法查字，必须熟记拼音字母的顺序，记得越熟，就越得越快。

为了查字又快又准，我们还可以直接去查字音的音节，然后按这个音节注明的页数很快查到要查的字。例如我们要
\newpage
\noindent
查“蝎”（jǐ）字，可以在《新华字典》前面所附的《汉语拼音音节索引》中，找到“j”起头的“ji”这个音节，根据这个音节标明的页码“218”，立即查到这个字；或者在《现代汉语词典》前面所附的《音节表》中，找到“j”起头的“ji”这个音节，根据这个音节标明的页码“568”，立即查到这个字。

2. 部首检字法

音序检字法比较科学，但是它有一个缺点，就是必须知道所查的那个汉字在普通话里怎样读，它是用哪些字母拼成的，如果自己不会读或者读不准，查起来就有困难了。因此，根据音序排列的字典或词典，一般都要和其它检字法配合着使用。例如《新华字典》、《现代汉语词典》等都附有《部首检字表》或者《四角号码检字表》。凡是自已不会读或者读不准的字，只能运用别的检字法来查。其中部首检字法是一种常用的检字法，应该掌握它。

什么是部首？就是字典、词典根据汉字形体偏旁所分的门类。人们按照字形结构，取其相同的偏旁部分，作为查字的依据。如：“姹、姥、娆、婆、娑”都有共同的偏旁部分“女”；“纶、绦、绐、缴、纍”都有共同的偏旁部分“糸”；“秀、香、秽、嵇、穑”都有共同的偏旁部分“禾”。这“女”、“糸”、“禾”，就是部首。

按部首排列的字典或词典，部首是按笔画多少的次序排列的，就是同一部首里的字，也是笔画少的排列在前，笔画多的排列在后。象《康熙字典》、《辞海》、《辞源》等，就是这样编排的。

使用这类字典或词典查字，关键在于要把所查的字的部
\newpage
\noindent
首和笔画搞清楚，才可能查得又快又准。因此，我们在使用这类字典或词典时，必须注意以下几点：

首先应知道要查的字典或词典里有哪些部首。如果能够熟悉字典或词典前边的《部首表》或者《部首目录》，并能熟悉部首的次序，那就更好。另外，还应注意各种字典的部首并不完全一样，例如《康熙字典》分二百一十四部，《新华字典》的《部首检字表》则对旧字典部首作了些调整和修改，只有一百八十九部。使用这些字典，必须知道它们分部的异同。

其次应知道要查的这个字在哪一部首里。哪一个字应该在哪一部里，一般比较容易看出来，例如“记”在“讠”（繁体“言”）部，“好”在“女”部，“地”在“土”部，有的就不大容易看出来，如“表、衰”在“衣”部，“邱、祁”在“邑”部，“防、陛”在“阜”部（《新华字典》新添了在右面的“阝”部和在左面的“阝”部，废除了“邑”部和“阜”部）；还有些字本身就是部首，如“穴”（不在“宀”部或“八”部），“音”（不在“立”部或“日”部），
\noindent
“麻”（不在“广”部或“木”部）等等。这类比较难查的字，我们可以留心地记一记，也可以在字典《部首表》后面所附的《检字表》或《难检字表》中按照字的笔画查出来。

另外还应知道要查的这个字的笔画数。通常说某字在某部多少画，是指部首以外的笔画数说的。例如“路”在“足（⻊）”部六画，“线”在“纟（糸）”部五画（繁体“綫”作八画），“暴”在“日”部十一画。另外，有些字往往会多数一笔或少数一笔，例如“仍”在“亻（人）”部二画，
\newpage
\noindent
“乃”在“丿”部一画，“长”在“丿”部三画（繁体“長”是部首），“依”在“亻”部四画（繁体“倶”是“亻”部八画），查字典或词典时，应将这类字的笔画数弄清并且记住，查起来就能又快又准了。

3. 笔画检字法

笔画检字法是按照字的笔画数和起笔的种类来检字的。许多字典或词典，如《中华大字典》、《同音字典》、《现代汉语词典》、《四角号码新词典》、《汉语成语小词典》、新《辞海》等，都附有《笔画检字表》或《笔画索引》，它也是目前使用较普遍的一种检字法，我们也应该掌握它。

这类笔画检字表所列各字，是按笔画多少排列的：少的在前，多的在后；笔画相同的，按几种起笔的次序排列；起笔相同的，再按第二笔的笔形次序排列，其余的以此类推。

以北京大学中文系一九五五级语言班编的《汉语成语小词典》为例，它把起笔分为七种，按“丶一丨丿㇀”次序排列。我们在查一个成语的第一个字（字头）时，就要以这七种起笔的顺序为依据。例如“群”字，可在十三画以“一”起笔的字下面查到；“漠”、“痴”两个字都在十三画，起笔都是“丶”，因为“漠”字第二笔是“一”，“痴”的第二笔是“丶”，所以“痴”一定要在“漠”的后面才能查到。这类笔画检字表，也有把起笔分为五种或四种的，在查字时，应注意有关说明。

另外，还要注意字典或词典所用铅字是宋体还是楷体，因为二者的笔画不尽相同。如“者”字宋体是九画，而楷体则为八画。还要注意某些字的笔画数。如“及”算三画还是四画（现在一般算三画），“巨”字算四画还是五画（现在
\newpage
\noindent
一般算四画）。在查字的时候，还要注意字的笔序，不按照一定的笔顺去查，就不容易查到。

4. 四角号码检字法

这种检字法是把每个字，不管什么形体，都分成四个角，把角的笔形分为十种，各用一个号码来表示。查字时，看字的四角的笔形是什么，依左上角、右上角、左下角、右下角的顺序，找出代表这个字的四角笔形的号码，然后按照这个号码在字典或词典里查出这个字来。由于这种检字法比较简便，使用得就比较广泛。例如《四角号码新词典》以及《新华字典》、《现代汉语词典》等所附的四角号码检字表，都是用的这种检字法，但具体的检字方法并不完全相同，后者在前者的基础上作了一些修改。
我们试以《现代汉语词典》所附的《四角号码检字表》为例。它把笔形分为十种，用0到9十个号码代表，并且列表如下：
\begin{tabular}{l|l|l|l|l|p{2cm}}
\hline
\multicolumn{2}{c|}{笔\quad名} & \multicolumn{1}{|c}{号码} & \multicolumn{1}{|c}{笔\quad形} & \multicolumn{1}{|c}{字\quad例} & \multicolumn{1}{|c}{说\quad明} \\
\hline
复笔 & 头 & 0 &  & 主病广言 & 点和横相结合 \\
\hline
单笔 & 横 & 1 & 一 & 天土 & 横 \\
\hline
单笔 & 横 & 1 & 丶丨丿 & 活络织见风 & 挑、横上钩和斜右钩 \\
\hline
单笔 & 垂 & 2 & 丨 & 旧山 & 直 \\
\hline
单笔 & 垂 & 2 & ノ丶 & 千顺力则 & 撇和直左钩 \\
\hline
单笔 &点 & 3 & 丶 & 宝社军外去亦 & 点 \\
\hline
单笔 &点 & 3 &  & 造瓜 & 捺 \\
\hline
\end{tabular}
\newpage
\noindent
（续表）

\noindent
\begin{tabular}{l|l|p{0.4cm}|p{1.6cm}|l|p{2.8cm}}
\hline
\multicolumn{2}{c|}{笔\quad名} & \multicolumn{1}{|c}{号码} & \multicolumn{1}{|c}{笔\quad形} & \multicolumn{1}{|c}{字\quad例} & \multicolumn{1}{|c}{说\quad明} \\
\hline
复笔 & 叉 & 4 & 十 & 古草 & 两笔交叉 \\
\hline
复笔 & 叉 & 4 & 才七又才 & 对式皮猪  & 两笔交叉 \\
\hline
复笔 & 串 & 5 & 丰 & 青本 & 一笔穿过两笔以上 \\
\hline
复笔 & 串 & 5 & 才戈丰 & 打戈泰申史 & 一笔穿过两笔以上 \\
\hline
复笔 & 方 & 6 & 口 & 另扣国甲由曲 & 四角整齐的方形 \\
\hline
复笔 & 方 & 6 & 口 & 目四	& 四角整齐的方形 \\
\hline
复笔 & 角 & 7 & 一ㄱㄴㄷ & 刀写亡表 & 一笔的转折 \\
\hline
复笔 & 角 & 7 & ㄱㄴㄷ & 阳兵又雪 & 两笔笔头相接所成的角形 \\
\hline
复笔 & 八 & 8 & 八 & 分共 & 八字形 \\
\hline
复笔 & 八 & 8 & 人入ㄱㄴㄷ & 余央籴羊午 & 八字形的变形 \\
\hline
复笔 & 小 & 9 & 小 & 尖宗 & 小字形 \\
\hline
复笔 & 小 & 9 & & 快木录当兴组 & 小字形的变形 \\
\hline
\end{tabular}
\vspace{1em}

表中所列的笔形和号码，必须记得烂熟，而且通过多次查字的实践，才可能查得又快又准。在查字之前，还必须细读表下所叙述的“查字方法”，以明确一般查法和特殊类型的查法。例如：

8 0 7 7 4 4 8 0
人 民 革 命

按照笔形得出: 8000, 7774, 4450, 8062，一查就可查到。其中的“人”、“革”、“命”三字，按照“查字方

\newpage
\noindent
法”。
法”的规定，“一个笔形，前角已经用过，后角作为0”，因此都有“0”的号码出现；甚至“人”字按笔形规定的号码是“8”，而其他各角都已用过这个笔形，所以其他各角都为“0”。又如四角同码字较多时，再取靠近右下角上方一个笔形作“附号”。如“民”字，附号的笔形为“丿”，作丁；“革”字，附号的笔形为“口”，作6；“命”字，附号的笔形为“一”，作7；“人”字没有附号的笔形，作0。因此，“人民革命”四字的号码为8000,7774,4450,8062。其它规定请参看《现代汉语词典》所载的《四角号码查字法》。

总之，要想又快又准地查字，必须首先熟悉各种检字法的规则，还得经过一番查字的实践才行。在具体运用一种字典或词典查字时，应在熟悉此字典、词典的检字法基础上，灵活配合其他检字法，以便迅速、准确地查到有些特殊的字。照此办理，自能越查越熟，从而达到得心应手、运用自如的地步。
